
% =============================================================================
% Project : THE PHILOSOPHY OF VIABILITY
% Book    : Book I — Ontology of Viability
% Book ID : VS.B01
% Part    : Part I — Identity and Manifestation
% Part ID : VS.B01.P01
% Chapter : Chapter 08 — Comparative Philosophical Examination
% Chapter ID : VS.B01.P01.C08
%
% File :
%   2_MAINMATTER/PART_I_IDENTITY_AND_MANIFESTATION/
%   CHAPTER_08_COMPARATIVE_PHILOSOPHICAL_EXAMINATION/
%   sections/section_16_ibn_arabi.tex
%
% Section : 16 — Ibn Arabi
% Section ID : VS.B01.P01.C08.S16
%
% Purpose :
%   A critical historical and ontological reconstruction of
%   Ibn Arabi's account of Being, divine unity, determination,
%   the divine names, self-disclosure, fixed entities,
%   receptivity, and the Barzakh, followed by a bilateral
%   comparison with the foundational commitments of Book I.
%
% =============================================================================

\section{Ibn Arabi}
\label{sec:ch08_s16}


% =============================================================================
% SS01 — SCOPE AND COMPARATIVE QUESTION
% =============================================================================

\subsection{Scope and Comparative Question}
\label{subsec:ch08_s16_ss01}

% Structural ID: VS.B01.P01.C08.S16.SS01.P001
\phantomsection
\label{para:VS.B01.P01.C08.S16.SS01.P001}

The comparison with Ibn Arabi introduces a philosophical
architecture in which the unity of the Real, the multiplicity
of determinate realities, and the manifestation of existence
are examined through a theological and metaphysical account
of divine self-disclosure. Its importance for the present
inquiry lies in the depth with which it addresses the relation
between what is unlimited and the differentiated forms
through which reality becomes manifest. The central question
is not whether Ibn Arabi and Book I use similar vocabulary.
It is whether their respective accounts of Being, Identity,
determination, and manifestation answer the same ontological
problem, and, where they do, whether they provide the same
explanation.

% Structural ID: VS.B01.P01.C08.S16.SS01.P002
\phantomsection
\label{para:VS.B01.P01.C08.S16.SS01.P002}

The comparison concentrates on the ontology of the Real
(al-Haqq), the meaning of existence (wujud), the relation
between unity and multiplicity, the divine names, tajalli,
the fixed entities (a'yan thabita), preparedness and
receptivity, and the Barzakh. These subjects are selected
because they bear directly upon the ontological relation
between unity, determination, and manifestation. The
discussion of imagination, the Perfect Human Being, and
the later systematization of the Akbarian school is included
only insofar as it clarifies this relation. Their complete
theological, epistemological, anthropological, and cosmological
implications are not developed here.

% Structural ID: VS.B01.P01.C08.S16.SS01.P003
\phantomsection
\label{para:VS.B01.P01.C08.S16.SS01.P003}

The governing comparative distinction is between two
different questions. Ibn Arabi asks how the one Real can
be the source and ground of the multiplicity of determinate
realities without being divided or exhausted by them.
Book I asks how Being and Identity are ontologically related,
and how manifestation can be examined without identifying
Identity with the determinations or manifestations through
which an existent becomes manifest. The questions intersect,
but neither can be substituted for the other. The comparison
therefore tests conceptual proximity while preserving the
independence of the two explanatory projects.


% =============================================================================
% SS02 — HISTORICAL AND TEXTUAL LOCATION
% =============================================================================

\subsection{Historical and Textual Location}
\label{subsec:ch08_s16_ss02}

% Structural ID: VS.B01.P01.C08.S16.SS02.P001
\phantomsection
\label{para:VS.B01.P01.C08.S16.SS02.P001}

Muhyi al-Din Ibn Arabi (1165--1240) was born in Murcia
and died in Damascus. His writings encompass theology,
Qur'anic interpretation, jurisprudence, metaphysics,
spiritual anthropology, and the interpretation of religious
experience. Two works are particularly important for the
present comparison: \textit{Fusus al-Hikam} (The Bezels
of Wisdom) and \textit{al-Futuhat al-Makkiyya} (The Meccan
Openings). The former presents a concentrated treatment
of prophetic wisdom through a sequence of chapters devoted
to particular prophetic figures; the latter contains a
much more extensive and varied exposition of his thought.
The two works should not be treated as interchangeable
summaries of a single, mechanically unified system
\citep{ChittickIbnArabiSEP2025}.

% Structural ID: VS.B01.P01.C08.S16.SS02.P002
\phantomsection
\label{para:VS.B01.P01.C08.S16.SS02.P002}

The distinction between Ibn Arabi's own formulations and
the subsequent Akbarian tradition is essential. Sadr al-Din
Qunawi and later interpreters played a substantial role in
systematizing, transmitting, and technically developing
concepts associated with Ibn Arabi. Consequently, a term
or scheme found in a later commentary cannot automatically
be attributed to Ibn Arabi in precisely the same form.
The Five Divine Presences, for example, acquired a
systematic role through later developments of the tradition.
Their use in reconstructing Ibn Arabi must therefore be
historically qualified \citep{ChittickFiveDivinePresences1982}.

% Structural ID: VS.B01.P01.C08.S16.SS02.P003
\phantomsection
\label{para:VS.B01.P01.C08.S16.SS02.P003}

The expression \textit{wahdat al-wujud}, commonly translated
as the Unity of Being or Oneness of Existence, requires
particular care. Ibn Arabi does not employ this expression
as the technical designation of his own doctrine. Its
subsequent attribution to him belongs to a history of
interpretation, controversy, and systematization. The
historical question must therefore be divided into two:
what Ibn Arabi actually says about the unity of the Real
and the existence of things, and how later authors classified
that position under the expression \textit{wahdat al-wujud}.
The distinction is indispensable for avoiding the projection
of later controversies onto the original texts
\citep{ChittickIbnArabiSEP2025,ChittickWahdatAlWujud1994}.


% =============================================================================
% SS03 — THE REAL AND THE MEANING OF WUJUD
% =============================================================================

\subsection{The Real and the Meaning of Wujud}
\label{subsec:ch08_s16_ss03}

% Structural ID: VS.B01.P01.C08.S16.SS03.P001
\phantomsection
\label{para:VS.B01.P01.C08.S16.SS03.P001}

The concept of \textit{wujud} occupies a central position
in Ibn Arabi's metaphysical vocabulary. It is commonly
translated as existence or Being, but its semantic range
also carries associations with finding, presence, and
encounter. Its use must be understood within the interaction
between philosophical discourse and the language of
theological and Sufi realization. Ibn Arabi draws upon
philosophical discussions of necessity and possibility
while placing the reality of existence within a broader
account of divine presence and manifestation. The concept
therefore cannot be reduced either to an abstract logical
predicate or to a description of a subject's psychological
experience \citep{ChittickSufiPathKnowledge1989}.

% Structural ID: VS.B01.P01.C08.S16.SS03.P002
\phantomsection
\label{para:VS.B01.P01.C08.S16.SS03.P002}

The Real, designated by the name \textit{al-Haqq}, occupies
the position of necessary and ultimate reality. Everything
other than the Real is considered in relation to the
possibility, contingency, and limitation of its existence.
The distinction between the Real and the possible is not
merely a distinction between two objects within a common
class. It concerns the difference between that whose
existence is necessary and that whose existence is dependent
upon the Real. The multiplicity of possible things therefore
does not introduce an independent source of existence
alongside the Real.

% Structural ID: VS.B01.P01.C08.S16.SS03.P003
\phantomsection
\label{para:VS.B01.P01.C08.S16.SS03.P003}

This account gives the unity of the Real a constitutive
role in the explanation of multiplicity. The many things
do not multiply the Real as though it were a material
whole divided into separate portions. Their differentiation
concerns the determinate modes in which existence becomes
manifest. The unity at issue is therefore not the numerical
unity of an aggregate, nor the simplicity of an ordinary
object. It concerns the non-duplication of the ultimate
reality of existence. This distinction is important because
the multiplicity of determinate realities remains a
substantive feature of the account rather than an illusion
to be dismissed.

% Structural ID: VS.B01.P01.C08.S16.SS03.P004
\phantomsection
\label{para:VS.B01.P01.C08.S16.SS03.P004}

The relation between the Real and the possible also
requires a distinction between existence in itself and
existence under limitation. The Real is discussed as
unrestricted or non-delimited existence
(\textit{al-wujud al-mutlaq}), whereas the cosmos is
considered through determinate or delimited existence
(\textit{al-wujud al-muqayyad}). This vocabulary describes
the relation between the unlimited and the differentiated.
It does not establish that the Real is a universal genus
of which individual things are species, nor that the
cosmos is a collection of independent pieces of the
divine Essence. The relation is one of dependence and
self-disclosure within a theological ontology
\citep[3:162.23]{IbnArabiFutuhat1911}.


% =============================================================================
% SS04 — TAWHID AND THE UNITY OF EXISTENCE
% =============================================================================

\subsection{Tawhid and the Unity of Existence}
\label{subsec:ch08_s16_ss04}

% Structural ID: VS.B01.P01.C08.S16.SS04.P001
\phantomsection
\label{para:VS.B01.P01.C08.S16.SS04.P001}

The principle of \textit{tawhid}, the affirmation of divine
unity, governs Ibn Arabi's treatment of existence. Its
metaphysical significance concerns the recognition that
the Real is not one existent among other independently
self-sufficient existents. The existence of the cosmos
is dependent, delimited, and related to the Real. This
dependence does not erase the distinctions among things;
rather, it determines the sense in which their existence
can be affirmed without assigning them an independent
source of reality.

% Structural ID: VS.B01.P01.C08.S16.SS04.P002
\phantomsection
\label{para:VS.B01.P01.C08.SS04.P002}

The later expression \textit{wahdat al-wujud} has been
used in different ways, including as a description of
the unity of ultimate existence, as a claim about the
relation between God and the world, and as a polemical
designation of positions alleged to collapse their
distinction. These interpretations are not equivalent.
In particular, the accusation that Ibn Arabi simply
identifies God with the aggregate of created things
does not adequately reconstruct the distinction between
the Real and the possible, or the distinction between
the Real in itself and its determinate manifestations.
The historical reception of the expression must be
examined separately from the metaphysical claims found
in the primary texts.

% Structural ID: VS.B01.P01.C08.S16.SS04.P003
\phantomsection
\label{para:VS.B01.P01.C08.SS04.P003}

The strongest form of the position is therefore not
that the world is numerically identical with God in
the ordinary sense, nor that the world is an independent
second reality. It is that the existence of the world
is intelligible only through its relation to the Real,
while the Real is not exhausted by the world or by
any particular determination. This formulation preserves
both dependence and differentiation. Whether it fully
explains the ontological status of each determinate
existent remains a central question for the comparison.


% =============================================================================
% SS05 — NONDELIMITATION AND TRANSCENDENCE
% =============================================================================

\subsection{Nondelimitation and Transcendence}
\label{subsec:ch08_s16_ss05}

% Structural ID: VS.B01.P01.C08.S16.SS05.P001
\phantomsection
\label{para:VS.B01.P01.C08.S16.SS05.P001}

Nondelimitation is not simply the absence of a boundary
within a spatial or quantitative field. It designates
the unrestricted character of the Real, whose reality
cannot be exhausted by any finite determination. The
Real is not one determinate form among others, because
every determinate form presupposes a limit or a respect
under which it is specified. The unlimited is therefore
not a larger object situated beyond the world; it is
the reality upon which the existence of determinate
things depends.

% Structural ID: VS.B01.P01.C08.S16.SS05.P002
\phantomsection
\label{para:VS.B01.P01.C08.SS05.P002}

The account of nondelimitation also prevents the Real
from being confined to transcendence understood as
absolute absence from the world. If the Real were
simply external to every possible manifestation, the
relation between the Real and the cosmos would remain
unexplained. Ibn Arabi's account therefore holds together
incomparability and disclosure: the Real is not identical
with any determinate form, yet no determinate form
possesses existence independently of the Real. The
language of transcendence and immanence must be read
within this relation rather than as two mutually
exclusive spatial locations.

% Structural ID: VS.B01.P01.C08.S16.SS05.P003
\phantomsection
\label{para:VS.B01.P01.C08.SS05.P003}

Nondelimitation must not be confused with invariance
in the technical sense developed in Book I. Invariance
requires an identified transformation, domain, respect,
type of variation, and limit. Nondelimitation concerns
the unrestricted reality of the Real and its relation
to determinate existence. It does not provide, by
itself, a criterion for identifying what remains
invariant through a transformation. Nor does it
establish that every manifestation preserves an
identical configuration. The two concepts address
different ontological questions.


% =============================================================================
% SS06 — THE DIVINE NAMES AND RELATIONAL DETERMINATION
% =============================================================================

\subsection{The Divine Names and Relational Determination}
\label{subsec:ch08_s16_ss06}

% Structural ID: VS.B01.P01.C08.S16.SS06.P001
\phantomsection
\label{para:VS.B01.P01.C08.S16.SS06.P001}

The divine names provide a central explanation of how
the unity of the Real is related to multiplicity.
Names such as the Living, the Knowing, the Merciful,
and the Powerful designate distinct meanings through
which the Real is spoken of and through which the
effects of divine attributes become manifest. Their
plurality does not imply a plurality of independent
divine essences. Rather, the names articulate different
relations and meanings associated with the Real.
Ibn Arabi's account thus gives relational differentiation
a central role in explaining the plurality of the
manifested world \citep{ChittickSufiPathKnowledge1989}.

% Structural ID: VS.B01.P01.C08.S16.SS06.P002
\phantomsection
\label{para:VS.B01.P01.C08.SS06.P002}

The names are not merely verbal labels imposed by
human observers. Their theological and metaphysical
function is not exhausted by the linguistic expressions
through which they are articulated. The distinction
between the names as spoken or written and the divine
realities to which they refer is essential. A linguistic
name belongs to a mode of expression; the divine name
concerns a relation or meaning within the account of
the Real and its manifestations. Consequently, the
multiplicity of names cannot be reduced to a plurality
of human descriptions.

% Structural ID: VS.B01.P01.C08.S16.SS06.P003
\phantomsection
\label{para:VS.B01.P01.C08.SS06.P003}

The names also clarify why the multiplicity of
manifestations does not divide the Real. Different
names become relevant through different determinate
relations and loci of manifestation. The same ultimate
reality is considered under different relational
meanings, while the manifested effects differ according
to the realities and capacities of the loci. This
structure provides a sophisticated account of
differentiation without requiring several independent
sources of existence. It does not, however, establish
that divine names are equivalent to the determinations
or relational structures discussed in Book I. The
theological status of the names remains specific to
Ibn Arabi's ontology.


% =============================================================================
% SS07 — TAJALLI AND SELF-DISCLOSURE
% =============================================================================

\subsection{Tajalli and Self-Disclosure}
\label{subsec:ch08_s16_ss07}

% Structural ID: VS.B01.P01.C08.S16.SS07.P001
\phantomsection
\label{para:VS.B01.P01.C08.S16.SS07.P001}

The concept of \textit{tajalli} is indispensable for
understanding the relation between the Real and the
manifested cosmos. It is commonly translated as
self-disclosure, self-manifestation, or theophany.
The concept concerns the disclosure of the Real through
determinate forms and relations. Manifestation is
therefore not merely the external appearance of an
already constituted object. It is a central feature
of the relation through which the existence and
attributes of the Real become manifest in the cosmos.

% Structural ID: VS.B01.P01.C08.S16.SS07.P002
\phantomsection
\label{para:VS.B01.P01.C08.SS07.P002}

The disclosure of the Real does not mean that the
divine Essence is exhausted or comprehensively
represented by any particular manifestation. Every
determinate form discloses the Real under a specific
respect while also concealing what cannot be contained
within that determination. The same manifestation
can therefore be considered as a disclosure and as
a veil. This double character is not a logical defect;
it follows from the relation between unrestricted
existence and the limited receptivity of determinate
forms. Tajalli thus explains both manifestation and
the incompleteness of every particular disclosure
\citep{ChittickSelfDisclosureGod1998}.

% Structural ID: VS.B01.P01.C08.S16.SS07.P003
\phantomsection
\label{para:VS.B01.P01.C08.SS07.P003}

The concept must not be reduced to the act of an
ordinary personal agent intentionally communicating
information to an observer. Its theological meaning
is more extensive: it concerns the ontological
self-disclosure of the Real and the dependence of
manifested things upon that disclosure. This does
not eliminate the personal and revelatory dimensions
of Ibn Arabi's thought; it prevents them from being
mistaken for the complete meaning of tajalli.
The concept therefore cannot be translated without
remainder into the generic notion of manifestation
used in Book I.

% Structural ID: VS.B01.P01.C08.S16.SS07.P004
\phantomsection
\label{para:VS.B01.P01.C08.SS07.P004}

The relation between tajalli and manifestation
constitutes one of the strongest points of comparison.
Both accounts assign ontological significance to
manifestation and reject its reduction to a merely
subjective appearance. Their difference concerns
the explanatory source and structure of manifestation.
In Ibn Arabi, manifestation is inseparable from the
self-disclosure of the Real. In Book I, manifestation
is examined as the manner or mode in which an existent
or an aspect of it is manifest. The latter formulation
does not presuppose a theological source or a
self-disclosing agent. The difference is foundational,
not merely terminological.


% =============================================================================
% SS08 — THE FIXED ENTITIES
% =============================================================================

\subsection{The Fixed Entities and Divine Knowledge}
\label{subsec:ch08_s16_ss08}

% Structural ID: VS.B01.P01.C08.S16.SS08.P001
\phantomsection
\label{para:VS.B01.P01.C08.S16.SS08.P001}

The doctrine of the fixed entities
(\textit{a'yan thabita}) addresses the status of
determinate realities in relation to divine knowledge.
The entities are known to the Real without possessing
independent existence in themselves. Their fixity
concerns their determinate realities as known by
God; it does not mean that they are already existing
objects located in a separate metaphysical region.
The doctrine therefore provides an account of how
the specificity of possible things can be related
to divine knowledge without assigning them an
independent source of existence
\citep{ChittickSufiPathKnowledge1989}.

% Structural ID: VS.B01.P01.C08.S16.SS08.P002
\phantomsection
\label{para:VS.B01.P01.C08.SS08.P002}

The distinction between fixed entities and existent
entities (\textit{a'yan mawjuda}) is not a distinction
between two unrelated kinds of things. It concerns
the same determinate reality considered under
different ontological conditions: as known in divine
knowledge and as manifested through the engendering
command. The entity does not acquire an independent
source of existence merely by becoming manifest.
Its existence remains dependent upon the Real, while
its specificity is not treated as an arbitrary
addition imposed by a finite observer. This relation
is central to Ibn Arabi's account of determination
and manifestation.

% Structural ID: VS.B01.P01.C08.S16.SS08.P003
\phantomsection
\label{para:VS.B01.P01.C08.SS08.P003}

The fixed entities must not be identified with Platonic
Forms without qualification. Nor should they be
translated simply as immutable essences if that
translation obscures the distinction between an
entity, its whatness, and the divine Essence. The
term \textit{'ayn} has a technical role that cannot
be replaced by a single familiar concept from another
metaphysical tradition. The fixed entity is not an
independently existing archetype that produces a
second object by imitation. It is the determinate
reality of a possible thing considered in relation
to divine knowledge and its subsequent manifestation
\citep{ChittickIbnArabiSEP2025}.

% Structural ID: VS.B01.P01.C08.S16.SS08.P004
\phantomsection
\label{para:VS.B01.P01.C08.S16.SS08.P004}

The relation between the fixed entity and Identity
in Book I is therefore a central point of divergence.
Identity in Book I is the ontological fact that an
existent is the existent it is. It is not defined
as an object of divine or human knowledge, a
pre-existing archetype, or a collection of determinate
properties. The fixed entity, by contrast, is
explained through its status in divine knowledge
and its relation to the possibility and manifestation
of existence. The two concepts address related
questions of specificity and determination, but
they do not possess the same explanatory function.
Equating them would replace an ontological question
about Identity with a specific theological account
of determinate realities.


% =============================================================================
% SS09 — PREPAREDNESS AND RECEPTIVITY
% =============================================================================

\subsection{Preparedness, Receptivity, and the Loci of Manifestation}
\label{subsec:ch08_s16_ss09}

% Structural ID: VS.B01.P01.C08.S16.SS09.P001
\phantomsection
\label{para:VS.B01.P01.C08.SS09.P001}

The manifestation of the divine names is not uniform
in every locus. Ibn Arabi employs the concepts of
preparedness (\textit{isti'dad}) and receptivity
(\textit{qabiliyya}) to explain why determinate
realities receive and display the effects of the
divine names in different ways. The Real's giving
is not restricted by a deficiency in the source;
the differentiated manifestation concerns the
specific capacities of the loci. Consequently,
multiplicity arises not from a division of the
Real but from the different ways in which its
disclosure becomes manifest \citep{ChittickSufiPathKnowledge1989}.

% Structural ID: VS.B01.P01.C08.S16.SS09.P002
\phantomsection
\label{para:VS.B01.P01.C08.SS09.P002}

Preparedness must not be interpreted as a causal
power independent of the Real. Nor is receptivity
merely a passive container into which an external
content is inserted. The concepts describe the
determinate relation between the reality of a thing,
the mode in which it can receive manifestation,
and the specific effects that become apparent in
it. This account gives the receiving locus a
differentiated role without making it an autonomous
source of existence.

% Structural ID: VS.B01.P01.C08.S16.SS09.P003
\phantomsection
\label{para:VS.B01.P01.C08.S16.SS09.P003}

The comparison with Book I must preserve the
difference between preparedness and ontological
Identity. A condition that explains how a thing
receives a manifestation does not thereby explain
what makes that existent the existent it is.
Likewise, a determinate capacity for manifestation
cannot be equated with Identity, Determination,
Structure, or Relational Determination in Book I.
The concepts may be compared with conditions and
respects relevant to manifestation, but their
theological grounding and their function within
Ibn Arabi's account remain distinct.


% =============================================================================
% SS10 — BARZAKH AND IMAGINAL ONTOLOGY
% =============================================================================

\subsection{The Barzakh and the Ontological Status of the Imaginal}
\label{subsec:ch08_s16_ss10}

% Structural ID: VS.B01.P01.C08.S16.SS10.P001
\phantomsection
\label{para:VS.B01.P01.C08.SS10.P001}

The concept of \textit{barzakh} is among Ibn Arabi's
most important contributions to the analysis of
intermediate realities. In its general sense, it
designates a limit or isthmus that both separates
and joins two sides. It is not simply a third object
positioned between two others. Its function is
relational: it establishes a boundary while also
making a relation between what it distinguishes
possible. The concept therefore provides an
ontological account of intermediacy that cannot
be reduced to spatial location or to a merely
conceptual division \citep{BashierIbnArabiBarzakh2004}.

% Structural ID: VS.B01.P01.C08.S16.SS10.P002
\phantomsection
\label{para:VS.B01.P01.C08.SS10.P002}

The Barzakh is employed in different contexts and
at different levels of Ibn Arabi's exposition.
It can designate the relation between existence
and nonexistence, the created cosmos in relation
to the Real, and the intermediate character of
the imaginal domain. Its meaning must therefore
be determined from the specific context in which
it appears. Treating every occurrence as the name
of one single cosmological region would obscure
the conceptual flexibility of the term.

% Structural ID: VS.B01.P01.C08.S16.SS10.P003
\phantomsection
\label{para:VS.B01.P01.C08.SS10.P003}

The imaginal realm (\textit{'alam al-khayal})
occupies a significant position in this account.
It is not merely a product of subjective fantasy
or an image constructed by the human mind. It
designates a mode or domain of reality in which
forms can possess an intermediate status between
the spiritual and the corporeal. The imaginal
therefore has ontological significance independent
of its recognition by a finite observer. Its
importance lies in showing that the structure
of reality cannot be exhausted by a simple
opposition between intelligible and sensible
existence \citep{ChittickSufiPathKnowledge1989}.

% Structural ID: VS.B01.P01.C08.S16.SS10.P004
\phantomsection
\label{para:VS.B01.P01.C08.S16.SS10.P004}

The Barzakh is relevant to Book I because it
demonstrates how an intermediate ontological
condition may possess a function that is neither
reducible to one side nor exhausted by the other.
Nevertheless, it is not equivalent to any of the
three mediating hubs in the architecture of Book I.
The latter have distinct proposed ontological
statuses and functions. Their relation to one
another cannot be derived from the general
meaning of the Barzakh. Nor can the imaginal realm
be automatically classified as one of the seven
realms of Book I. Such a classification would
require an independent argument and an explicit
decision concerning the architecture of Part II.


% =============================================================================
% SS11 — UNITY, MULTIPLICITY, AND THE ONE/MANY
% =============================================================================

\subsection{Unity, Multiplicity, and the One and the Many}
\label{subsec:ch08_s16_ss11}

% Structural ID: VS.B01.P01.C08.S16.SS11.P001
\phantomsection
\label{para:VS.B01.P01.C08.S16.SS11.P001}

The relation between unity and multiplicity is
not resolved by eliminating one of its terms.
The unity of the Real remains fundamental, while
the multiplicity of names, determinate realities,
and manifestations is also affirmed. The One is
not multiplied into independent ultimate sources
by the appearance of many determinate things.
At the same time, the many are not simply erased
from the account. Their reality is understood
through their dependence upon the Real and their
specific modes of manifestation.

% Structural ID: VS.B01.P01.C08.S16.SS11.P002
\phantomsection
\label{para:VS.B01.P01.C08.SS11.P002}

This account must be distinguished from a simple
doctrine of aggregate unity. The Real is not a
whole composed of its manifestations as parts,
and the plurality of manifestations does not
operate as a numerical division of the divine
Essence. The distinction between the unity of
the Real and the multiplicity of its names and
manifestations is therefore not the same as the
distinction between a whole and its components.
The metaphysical relation is grounded in the
dependence of determinate existence upon the
unrestricted Real.

% Structural ID: VS.B01.P01.C08.S16.SS11.P003
\phantomsection
\label{para:VS.B01.P01.C08.S16.SS11.P003}

For Book I, the decisive question concerns the
level at which unity and multiplicity are asserted.
Unity does not automatically establish numerical
singularity, and plurality among manifestations
does not establish a plurality of identities.
The possibility of multiple identities remains
an open ontological problem. Ibn Arabi's account
offers a powerful model of unity with differentiated
manifestation, but it does not settle the independent
question of the ontological status and possible
plurality of Identity as formulated in Book I.


% =============================================================================
% SS12 — RENEWAL AND NON-REPETITION OF TAJALLI
% =============================================================================

\subsection{Renewal and the Non-Repetition of Self-Disclosure}
\label{subsec:ch08_s16_ss12}

% Structural ID: VS.B01.P01.C08.S16.SS12.P001
\phantomsection
\label{para:VS.B01.P01.C08.S16.SS12.P001}

Ibn Arabi's account of continuous renewal gives
a dynamic dimension to the relation between the
Real and the cosmos. The doctrine of the renewal
of creation at each instant
(\textit{tajdid al-khalq fi'l-anat}) is connected
with the principle that self-disclosure is not
repeated in precisely the same manner. The
manifested world is therefore not a static
collection of forms whose existence remains
unchanged through time. Its determinate forms
are continuously renewed through the relation
between the Real and the loci of manifestation
\citep{ChittickSelfDisclosureGod1998}.

% Structural ID: VS.B01.P01.C08.S16.SS12.P002
\phantomsection
\label{para:VS.B01.P01.C08.SS12.P002}

The non-repetition of tajalli must not be interpreted
as proof that every manifestation belongs to a
different identity. It concerns the renewed
character of self-disclosure and the absence of
a simple repetition of the same manifestation.
A conclusion about Identity would require an
additional argument. Likewise, the renewal of
manifested forms cannot be equated automatically
with the destruction and replacement of the
identity of an existent. The relation between
manifestational renewal and ontological continuity
must remain analytically open.

% Structural ID: VS.B01.P01.C08.S16.SS12.P003
\phantomsection
\label{para:VS.B01.P01.C08.SS12.P003}

This doctrine raises a serious comparative problem
for any account that seeks to explain persistence
through the recurrence of the same manifested
configuration. If manifestations are renewed
without exact repetition, continuity cannot be
explained simply by the recurrence of identical
appearances. Book I must therefore distinguish
the question of what is manifested from the
question of whether and how Identity obtains
through transformations. The Akbarian account
does not supply a ready-made answer to that
question, but it makes the inadequacy of a
manifestation-only criterion especially clear.


% =============================================================================
% SS13 — THE PERFECT HUMAN AND LATER SYSTEMATIZATION
% =============================================================================

\subsection{The Perfect Human and the Comprehensive Structure of Reality}
\label{subsec:ch08_s16_ss13}

% Structural ID: VS.B01.P01.C08.S16.SS13.P001
\phantomsection
\label{para:VS.B01.P01.C08.S16.SS13.P001}

The concept of the Perfect Human
(\textit{al-insan al-kamil}) expresses the
comprehensive relation between the divine and
cosmic dimensions of Ibn Arabi's account. The
human being is considered in relation to the
divine names and to the differentiated structure
of the cosmos. The Perfect Human is not simply
an exemplary individual possessing a collection
of superior qualities. The concept concerns
the comprehensive capacity through which the
different dimensions of reality can be brought
into a unified manifestation. Its significance
is therefore both anthropological and ontological
\citep{ChittickSelfDisclosureGod1998}.

% Structural ID: VS.B01.P01.C08.S16.SS13.P002
\phantomsection
\label{para:VS.B01.P01.C08.SS13.P002}

The concept is closely related to the idea of
a comprehensive reality that mediates between
the Real and the cosmos. Later Akbarian authors,
particularly Qunawi and his successors, developed
systematic accounts of this mediating and
comprehensive function. Their formulations
must be distinguished from the precise wording
and organization of Ibn Arabi's own texts.
A comparison that treats every later scheme
as a direct statement by Ibn Arabi would
confuse the original writings with their
subsequent philosophical elaboration.

% Structural ID: VS.B01.P01.C08.S16.SS13.P003
\phantomsection
\label{para:VS.B01.P01.C08.SS13.P003}

The Perfect Human does not have a direct
counterpart in the foundational architecture
of Book I. The human being is not assigned
the role of an ontological center merely because
human beings can understand, interpret, or
reconstruct manifestations. Such a move would
make the existence and structure of reality
dependent upon the human knower. The comparative
value of Ibn Arabi's account lies instead in
its demonstration that a comprehensive relation
between different dimensions of reality can
be articulated without reducing them to a
single undifferentiated level. The specific
theological and anthropological role of the
Perfect Human remains distinct.


% =============================================================================
% SS14 — THE STRONGEST IBN ARABI CHALLENGE TO BOOK I
% =============================================================================

\subsection{The Strongest Ibn Arabi Challenge to Book I}
\label{subsec:ch08_s16_ss14}

% Structural ID: VS.B01.P01.C08.S16.SS14.P001
\phantomsection
\label{para:VS.B01.P01.C08.SS14.P001}

The strongest challenge that Ibn Arabi's account
poses to Book I concerns the explanatory necessity
of Identity as an ontological concept analytically
distinguished from Being. If all determinate
existence is dependent upon the Real, and if the
specificity of things is articulated through
their realities in divine knowledge, the question
arises whether an additional concept of Identity
performs an independent explanatory function.
The distinction between Being and Identity may
appear unnecessary if the account of the Real,
the divine names, and the fixed entities already
explains both existence and the specificity
of manifested things.

% Structural ID: VS.B01.P01.C08.S16.SS14.P002
\phantomsection
\label{para:VS.B01.P01.C08.SS14.P002}

A second challenge concerns the status of
determination. Ibn Arabi's account relates the
specificity of things to their realities and
to the divine knowledge of those realities.
Book I must explain whether Identity adds
anything beyond the determinations through
which an existent is examined. If Identity
were merely another name for a determinate
essence, a fixed entity, or a set of properties,
the distinction would lack independent
explanatory force. The challenge is therefore
not merely theological; it concerns the
conceptual adequacy of the proposed ontology.

% Structural ID: VS.B01.P01.C08.S16.SS14.P003
\phantomsection
\label{para:VS.B01.P01.C08.SS14.P003}

A third challenge concerns unity. The affirmation
of the unity of the Real raises the question
whether a plurality of identities is compatible
with the ontological dependence of all possible
existents upon one ultimate reality. Book I
cannot answer this by asserting that plurality
of manifestations entails plurality of identities.
Such an inference would violate the distinction
between manifestation and Identity. Nor can
the question be settled by treating the unity
of the Real as equivalent to the unity of
Identity. The relevant concepts must first
be distinguished before their relation can
be established.

% Structural ID: VS.B01.P01.C08.S16.SS14.P004
\phantomsection
\label{para:VS.B01.P01.C08.SS14.P004}

The challenge can be stated in its strongest
form: if Being and Identity are inseparable,
but neither is reducible to the other, what
ontological work does the analytical distinction
perform? What prevents Identity from becoming
a second explanatory principle added to Being?
And how can its irreducibility be established
without reducing it to essence, determination,
or the known specificity of a possible existent?
These questions cannot be answered by rejecting
the Akbarian account in advance. They require
Book I to articulate the positive content of
Identity with greater precision.


% =============================================================================
% SS15 — THE BOOK I CHALLENGE AND A POSSIBLE AKBARIAN REPLY
% =============================================================================

\subsection{The Book I Challenge and a Possible Akbarian Reply}
\label{subsec:ch08_s16_ss15}

% Structural ID: VS.B01.P01.C08.S16.SS15.P001
\phantomsection
\label{para:VS.B01.P01.C08.S16.SS15.P001}

The strongest challenge from Book I concerns
the distinction between the ultimate source
of existence and the ontological specificity
of an existent. Ibn Arabi explains the existence
of determinate things through their relation
to the Real, their realities in divine knowledge,
and their modes of manifestation. Book I asks
whether this account fully distinguishes
the fact that an existent is from the fact
that it is the existent it is. The question
is not whether Ibn Arabi recognizes determinate
realities. He clearly does. It is whether
the relation between their determinate
specificity and their existence provides
an account of Identity in the sense required
by Book I.

% Structural ID: VS.B01.P01.C08.S16.SS15.P002
\phantomsection
\label{para:VS.B01.P01.C08.SS15.P002}

A second challenge concerns manifestation.
Tajalli is a theological and ontological account
of the Real's self-disclosure. Manifestation
in Book I is not defined through a necessary
relation to divine self-disclosure. If the
Akbarian account is treated as a universal
definition of manifestation, it risks making
a particular theological explanation
the condition of every possible ontological
manifestation. Book I therefore asks whether
the broader category of manifestation can
be articulated without presupposing the
specific theological commitments of tajalli.

% Structural ID: VS.B01.P01.C08.S16.SS15.P003
\phantomsection
\label{para:VS.B01.P01.C08.SS15.P003}

A third challenge concerns the Barzakh.
Its capacity to unite and distinguish two
sides provides an important model of
ontological intermediacy. Yet the concept
does not, by itself, determine the number,
status, or functions of the mediating
domains in Book I. The Barzakh may be
structurally comparable to a mediating
condition, but it cannot replace the
separate analysis required for each
proposed hub. The distinction is especially
important because the architecture of
Book I does not treat all intermediate
domains as equivalent or as a single
undifferentiated level.

% Structural ID: VS.B01.P01.C08.S16.SS15.P004
\phantomsection
\label{para:VS.B01.P01.C08.SS15.P004}

An informed Akbarian reply could argue that
these objections presuppose a conceptual
separation that the unity of the Real
renders unnecessary. From this perspective,
the specificity of a thing is already
accounted for through its reality, its
relation to the divine names, and the
particular manner in which it receives
self-disclosure. The reply would question
whether the Book I distinction adds
explanatory content or merely introduces
another vocabulary for a relation already
explained. This reply is philosophically
serious. It cannot be dismissed by asserting
that the two systems use different terms;
Book I must demonstrate the independent
function of its own distinctions.


% =============================================================================
% SS16 — BILATERAL EXPLANATORY COSTS AND SUBSTITUTION
% =============================================================================

\subsection{Bilateral Explanatory Costs and the Substitution Test}
\label{subsec:ch08_s16_ss16}

% Structural ID: VS.B01.P01.C08.S16.SS16.P001
\phantomsection
\label{para:VS.B01.P01.C08.S16.SS16.P001}

The Akbarian account gains considerable
explanatory power by grounding the multiplicity
of determinate realities in one ultimate
source of existence. It connects the divine
names, the specificity of things, and their
manifestation within a unified theological
ontology. Its explanatory cost, from the
perspective of Book I, is that the irreducible
ontological role of Identity is not formulated
as a separate analytical problem. The account
must therefore be examined to determine
whether the specificity of things can be
fully explained through their fixed realities,
divine relations, and modes of disclosure.

% Structural ID: VS.B01.P01.C08.S16.SS16.P002
\phantomsection
\label{para:VS.B01.P01.C08.S16.SS16.P002}

Book I gains a distinct analytical space
for examining Identity without reducing
it to Being, Determination, Unity, or
Manifestation. Its explanatory cost is
the obligation to show that this distinction
does not introduce an unnecessary ontological
multiplication. It must establish why
Identity cannot be replaced by a more
familiar category such as essence, quiddity,
or determinate reality. It must also explain
how its distinction between Being and Identity
remains consistent with their ontological
inseparability and the absence of ontological
priority.

% Structural ID: VS.B01.P01.C08.S16.SS16.P003
\phantomsection
\label{para:VS.B01.P01.C08.SS16.P003}

The substitution test clarifies the difference.
Replacing Identity in Book I with the fixed
entity of Ibn Arabi would import the specific
relation between divine knowledge, nonexistence
in itself, and manifestation through the
engendering command. Replacing Being with
the Real or with non-delimited existence
would import a theological account of
ultimate reality that Book I has not adopted.
Replacing manifestation with tajalli would
make self-disclosure a necessary feature
of every manifestation. Replacing the
mediating hubs with the Barzakh would
erase distinctions among their proposed
functions. None of these substitutions
is conceptually neutral. The resemblance
between the systems is therefore real
but limited, and their explanatory
architectures remain distinct.


% =============================================================================
% SS17 — IMPLICATIONS FOR BEING AND IDENTITY
% =============================================================================

\subsection{Implications for Being and Identity}
\label{subsec:ch08_s16_ss17}

% Structural ID: VS.B01.P01.C08.S16.SS17.P001
\phantomsection
\label{para:VS.B01.P01.C08.S16.SS17.P001}

The comparison establishes that a metaphysics
of unity and manifestation can preserve
a substantial account of multiplicity
without treating the many as independent
ultimate sources of existence. Ibn Arabi's
analysis of the Real, the divine names,
the fixed entities, and receptivity offers
a highly developed explanation of how
determinate forms can be related to an
unlimited source. Its significance for
Book I lies in the precision with which
it articulates the dependence of manifested
realities and the differentiated character
of their disclosure.

% Structural ID: VS.B01.P01.C08.S16.SS17.P002
\phantomsection
\label{para:VS.B01.P01.C08.S16.SS17.P002}

The comparison also clarifies why Identity
cannot be identified with any one of the
Akbarian concepts considered here. The
fixed entity addresses a determinate reality
in relation to divine knowledge; the divine
name articulates a relation or meaning
through which the Real is considered;
tajalli concerns self-disclosure; and
preparedness concerns the differentiated
receptivity of the loci. Each concept
performs a particular explanatory function.
None, without further argument, provides
the full ontological meaning of Identity
as formulated in Book I.

% Structural ID: VS.B01.P01.C08.S16.SS17.P003
\phantomsection
\label{para:VS.B01.P01.C08.SS17.P003}

The distinction between Being and Identity
must therefore remain analytically explicit.
Their ontological inseparability does not
make them synonymous, and their analytical
distinction does not establish two independent
entities. The comparison with Ibn Arabi
reinforces the need to distinguish the
unity of ultimate existence from the
question of what makes an existent the
existent it is. The latter cannot be
settled solely by describing the source
of existence or the conditions under
which a manifestation becomes determinate.

% Structural ID: VS.B01.P01.C08.S16.SS17.P004
\phantomsection
\label{para:VS.B01.P01.C08.S16.SS17.P004}

The following diagram presents a schematic
reconstruction of the principal relations
in Ibn Arabi's account and distinguishes
them from the foundational relation
between Being and Identity in Book I.
It does not represent the Akbarian ontology
as a single linear causal sequence, nor
does it identify the fixed entities with
Identity or the divine names with
Determination.

% =============================================================================
% Project : THE PHILOSOPHY OF VIABILITY
% Book    : Book I — Ontology of Viability
% Book ID : VS.B01
% Part    : Part I — Identity and Manifestation
% Part ID : VS.B01.P01
% Chapter : Chapter 08 — Comparative Philosophical Examination
% Chapter ID : VS.B01.P01.C08
%
% File :
%   4_ASSETS/diagrams/
%   VS.B01.P01.C08.S16.SS17.P004.DG01.ibn_arabi_and_book_i_relational_distinctions.tex
%
% Asset : DG01
% Structural Identifier :
%   VS.B01.P01.C08.S16.SS17.P004.DG01
%
% Purpose :
%   To represent the principal relations in Ibn Arabi's ontology
%   and distinguish them from the foundational relation between
%   Being and Identity in Book I.
%
% Philosophical Limits :
%   1. The diagram is a conceptual reconstruction, not a literal
%      reproduction of a cosmological diagram by Ibn Arabi.
%   2. Connections do not necessarily represent temporal or
%      efficient-causal sequences.
%   3. Fixed entities are not identified with Book I Identity.
%   4. Divine names are not identified with Determination.
%   5. Tajalli is not identified with generic Manifestation.
%   6. The diagram does not establish an equivalence between
%      the Barzakh and the three mediating hubs of Book I.
%
% Dependencies :
%   tikz, positioning, arrows.meta.
%   (preamble: \usetikzlibrary{positioning, arrows.meta})
%
% =============================================================================

\begin{figure}[htbp]
\centering

\begin{tikzpicture}[
    font=\small,
    every node/.style={inner sep=7pt},
    realnode/.style={
        draw=ViabilityDarkBlue, fill=ViabilityPaleBlue,
        rounded corners=3pt, line width=1pt,
        text width=10.5cm, align=center, minimum height=1.5cm
    },
    namesnode/.style={
        draw=ViabilityGold, fill=ViabilityPaleGold,
        rounded corners=2pt, line width=0.8pt,
        text width=3.5cm, align=center, minimum height=2.8cm
    },
    entitiesnode/.style={
        draw=ViabilityBlue, fill=ViabilityPaleBlue,
        rounded corners=2pt, line width=0.8pt,
        text width=3.5cm, align=center, minimum height=2.8cm
    },
    receptivitynode/.style={
        draw=ViabilityGreen, fill=ViabilityPaleGreen,
        rounded corners=2pt, line width=0.8pt,
        text width=3.5cm, align=center, minimum height=2.8cm
    },
    manifestationnode/.style={
        draw=ViabilityDarkBlue, fill=ViabilityPaper,
        rounded corners=3pt, line width=1pt,
        text width=10.5cm, align=center, minimum height=1.6cm
    },
    booknode/.style={
        draw=ViabilityDarkBlue, fill=ViabilityPaleBlue,
        rounded corners=2pt, line width=0.8pt,
        text width=3.2cm, align=center, minimum height=1.2cm
    },
    identitynode/.style={
        draw=ViabilityGold, fill=ViabilityPaleGold,
        rounded corners=2pt, line width=0.8pt,
        text width=3.2cm, align=center, minimum height=1.2cm
    },
    relation/.style={
        -{Stealth[length=2mm]},
        draw=ViabilityDarkBlue, line width=0.75pt
    },
    conceptual/.style={
        densely dashed, draw=ViabilityGray, line width=0.8pt
    },
    title/.style={
        font=\small\bfseries, align=center, text width=12cm
    },
    label/.style={
        font=\scriptsize, inner sep=2pt, fill=white
    },
    note/.style={
        draw=ViabilityLightGray, fill=white, rounded corners=2pt,
        text width=11cm, align=center
    }
]

% ---------------------------------------------------------------------
% PANEL A — IBN ARABI
% ---------------------------------------------------------------------

\node[realnode] (real) at (0,0) {
    \textbf{The Real — al-Haqq}\\[3pt]
    Non-delimited Existence (\textit{al-wujud al-mutlaq})
};

\node[title, above=2mm of real] (titleA)
    {Ibn Arabi: The Real and Its Determinate Manifestations};

\node[entitiesnode, below=1.6cm of real] (entities) {
    \textbf{Fixed Entities}\\[4pt]
    \textit{A'yan thabita}\\
    Determinate realities\\
    Known in divine knowledge
};

\node[namesnode, left=6mm of entities] (names) {
    \textbf{Divine Names}\\[4pt]
    Relational meanings\\
    Divine attributes\\
    Distinct aspects of disclosure
};

\node[receptivitynode, right=6mm of entities] (receptivity) {
    \textbf{Preparedness and Receptivity}\\[4pt]
    \textit{Isti'dad / Qabiliyya}\\
    Specific capacity of the receiving loci
};

\node[manifestationnode, below=1.4cm of entities] (tajalli) {
    \textbf{Tajalli — Self-Disclosure}\\[4pt]
    The manifestation of the Real under determinate relations
    and according to the receptivity of the loci
};

\node[manifestationnode, below=1.2cm of tajalli] (forms) {
    \textbf{Delimited Manifestations}\\[4pt]
    Existent entities (\textit{a'yan mawjuda})\\
    Differentiated forms of manifested existence
};

% Relations from the Real (vertical drops)

\draw[relation] (names.north |- real.south) -- (names.north)
    node[pos=0.5, left, label, align=right] {relational\\meanings};

\draw[relation] (entities.north |- real.south) -- (entities.north)
    node[pos=0.5, right, label] {divine knowledge};

\draw[relation] (receptivity.north |- real.south) -- (receptivity.north)
    node[pos=0.5, right, label, align=left] {relation to\\the loci};

% Relations to self-disclosure (dashed, non-causal)

\draw[conceptual] (names.south) -- (names.south |- tajalli.north);
\draw[conceptual] (entities.south) -- (tajalli.north);
\draw[conceptual] (receptivity.south) -- (receptivity.south |- tajalli.north);

\draw[relation] (tajalli) -- (forms)
    node[pos=0.5, right, label] {manifested forms};

% ---------------------------------------------------------------------
% PANEL B — BOOK I
% ---------------------------------------------------------------------

\node[title, below=1.5cm of forms] (titleB)
    {Book I: The Foundational Relation};

\node[booknode, below=5mm of titleB, xshift=-4.2cm] (being)
    {\textbf{Being}};

\node[identitynode, below=5mm of titleB, xshift=4.2cm] (identity)
    {\textbf{Identity}};

\draw[conceptual] (being.east) -- (identity.west)
    node[pos=0.5, label, text width=3.3cm, align=center] {
        Analytical Distinction\\
        Ontological Inseparability\\
        Absence of Ontological Priority
    };

\coordinate (axisB) at (titleB |- being.south);

\node[note, below=6mm of axisB] {
    \textbf{Non-equivalence}\\[3pt]
    Fixed Entity $\neq$ Identity \quad Divine Name $\neq$ Determination\\
    Tajalli $\neq$ Generic Manifestation \quad
    Barzakh $\neq$ Book I Mediating Hubs
};

\end{tikzpicture}

\caption[
    Ibn Arabi's ontology and the foundational relation
    between Being and Identity in Book I
]{
    Schematic reconstruction of the relations among
    the Real, divine names, fixed entities, receptivity,
    and self-disclosure in Ibn Arabi's ontology.
    The lower panel separately presents the foundational
    relation between Being and Identity in Book I.
    The diagram is comparative rather than genealogical
    or causal; it does not identify the corresponding
    concepts across the two systems. The connections
    indicate conceptual relations within the reconstructed
    account and do not assert a single temporal or
    efficient-causal sequence.
}

\label{fig:VS.B01.P01.C08.S16.SS17.P004.DG01}

\end{figure}

% =============================================================================
% END OF DIAGRAM DG01
% =============================================================================


% =============================================================================
% SS18 — OPEN PROBLEMS AND CLAIM STATUS
% =============================================================================

\subsection{Open Problems and Claim Status}
\label{subsec:ch08_s16_ss18}

% Structural ID: VS.B01.P01.C08.S16.SS18.P001
\phantomsection
\label{para:VS.B01.P01.C08.S16.SS18.P001}

Several questions remain unresolved. First,
does the distinction between the fixed entity
and the existent entity provide an account
of ontological specificity that can be
compared with Identity without reducing
Identity to divine knowledge? Second,
does the relation between the Real and
the determinate world fully explain the
ontological status of each existent, or
does it leave a further question concerning
the identity of the existent? Third, how
should the relation between the non-repetition
of tajalli and the continuity of determinate
existents be formulated? These questions
remain open because the comparison has
not established a complete translation
between the two systems.

% Structural ID: VS.B01.P01.C08.S16.SS18.P002
\phantomsection
\label{para:VS.B01.P01.C08.S16.SS18.P002}

Further questions concern the relation
between the Barzakh and the mediating
domains of Book I, the status of the
imaginal realm within a broader ontology,
and the relation between the divine names
and the multiplicity of determinate
realities. The place of Light in Ibn
Arabi's ontology also requires careful
comparison with the independent treatment
of Light in Part III. None of these
questions should be settled merely by
terminological resemblance. Each requires
a separate examination of the relevant
ontological commitments and explanatory
functions.

% Structural ID: VS.B01.P01.C08.S16.SS18.P003
\phantomsection
\label{para:VS.B01.P01.C08.SS18.P003}

The claim-status record distinguishes
the functional or logical form of a claim
from its epistemic or justificatory status.
The historical reconstruction of Ibn
Arabi's position is not itself a
foundational commitment of Book I.
Likewise, a comparative thesis does not
become an adopted principle merely because
it identifies a point of convergence.
The distinction between the two axes
prevents the historical account, the
interpretation of primary texts, and
the commitments of Book I from being
collapsed into a single classification.

% Structural ID: VS.B01.P01.C08.S16.SS18.P004
\phantomsection
\label{para:VS.B01.P01.C08.S16.SS18.P004}

The following table records the principal
claims, their logical form, their
justificatory status, and the limits
of their use in the present comparison.


% =============================================================================
% Project : THE PHILOSOPHY OF VIABILITY
% Book    : Book I — Ontology of Viability
% Book ID : VS.B01
% Part    : Part I — Identity and Manifestation
% Part ID : VS.B01.P01
% Chapter : Chapter 08 — Comparative Philosophical Examination
% Chapter ID : VS.B01.P01.C08
%
% File :
%   4_ASSETS/tables/
%   VS.B01.P01.C08.S16.SS18.P004.TL01.ibn_arabi_claims_and_comparative_limits.tex
%
% Asset : TL01
% Structural Identifier :
%   VS.B01.P01.C08.S16.SS18.P004.TL01
%
% Purpose :
%   To distinguish the reconstructed claims of Ibn Arabi
%   from their logical forms and epistemic statuses,
%   while recording their comparative limits for Book I.
%
% Dependencies :
%   longtable, booktabs, array.
%
% =============================================================================

\begingroup
\small
\renewcommand{\arraystretch}{1.2}
\setlength{\tabcolsep}{4pt}
\setlength{\LTcapwidth}{\linewidth}

\def\colw#1{%
    \dimexpr(\linewidth-8\tabcolsep)*#1/100\relax
}

\def\rowsep{%
    \specialrule{0.2pt}{3pt}{3pt}
}

\def\tablehead{%
    \toprule
    \textbf{Concept or Claim} &
    \textbf{Akbarian Reconstruction} &
    \textbf{Logical Form} &
    \textbf{Epistemic Status} &
    \textbf{Comparative Limit for Book I}
    \\
    \midrule
}

\def\grouprow#1{%
    \multicolumn{5}{@{}l}{\footnotesize\textsc{#1}}\\
    \cmidrule{1-5}
    \addlinespace[1pt]
}

\begin{longtable}{@{}
    >{\raggedright\arraybackslash}p{\colw{14}}
    >{\raggedright\arraybackslash}p{\colw{29}}
    >{\raggedright\arraybackslash}p{\colw{14}}
    >{\raggedright\arraybackslash}p{\colw{17}}
    >{\raggedright\arraybackslash}p{\colw{26}}
@{}}

\caption{
    Ibn Arabi's Principal Ontological Claims and Their
    Comparative Limits
}
\label{tab:VS.B01.P01.C08.S16.SS18.P004.TL01}
\\

\tablehead
\endfirsthead

\multicolumn{5}{@{}l}{%
    \footnotesize\textit{
        Table~\thetable\ (continued):
        Ibn Arabi's Principal Ontological Claims
    }
}
\\

\tablehead
\endhead

\multicolumn{5}{@{}r}{%
    \footnotesize\textit{Continued on next page}
}
\\
\endfoot

\bottomrule
\endlastfoot

\grouprow{Unity and the Real}

\textbf{Tawhid}
&
The affirmation of the unity of the Real and the
dependence of possible existence upon it.
&
Thesis
&
Argued/Derived Claims
&
The unity of the Real is not equivalent to the
unity or singularity of Identity.
\\
\rowsep

\textbf{Wahdat al-wujud}
&
A later expression used in the history of interpretation;
not Ibn Arabi's own technical designation.
&
Historical interpretation
&
Argued/Derived Claims
&
The later label must not be projected without qualification
onto the primary texts.
\\
\rowsep

\textbf{Nondelimitation}
&
The unrestricted character of the Real in relation
to determinate and delimited existence.
&
Definition (historical reconstruction)
&
Argued/Derived Claims
&
Nondelimitation is not equivalent to Invariance,
Immutability, or Identity.
\\
\rowsep

\grouprow{Names and manifestation}

\textbf{Divine names}
&
Relational meanings and attributes through which
the Real is considered under differentiated respects.
&
Thesis
&
Argued/Derived Claims
&
Divine names are not identical with Determination
in Book I.
\\
\rowsep

\textbf{Tajalli}
&
The self-disclosure of the Real through determinate
forms and relations.
&
Thesis
&
Argued/Derived Claims
&
Tajalli is not equivalent to generic Manifestation.
\\
\rowsep

\textbf{Preparedness}
&
The specific capacity of a locus to receive and
display determinate manifestations.
&
Thesis
&
Argued/Derived Claims
&
Preparedness is not Identity or a complete
account of the conditions of manifestation in Book I.
\\
\rowsep

\grouprow{Determinate realities}

\textbf{A'yan thabita}
&
Determinate realities considered in relation to
divine knowledge without independent existence
in themselves.
&
Thesis
&
Argued/Derived Claims
&
Fixed entities must not be identified with Identity
or with independently existing Platonic archetypes.
\\
\rowsep

\textbf{A'yan mawjuda}
&
The same determinate realities considered through
their manifested or delimited existence.
&
Thesis
&
Argued/Derived Claims
&
Their relation to fixed entities does not establish
a Book I theory of identity through transformation.
\\
\rowsep

\grouprow{Intermediate realities}

\textbf{Barzakh}
&
A limit or isthmus that simultaneously separates
and joins two sides.
&
Definition (historical reconstruction)
&
Argued/Derived Claims
&
The Barzakh is not equivalent to any one of
Book I's three mediating hubs.
\\
\rowsep

\textbf{Imaginal realm}
&
An ontologically significant intermediate domain
between spiritual and corporeal forms.
&
Thesis
&
Conditionally Supported Claims
&
It is not automatically a separate realm
within Book I's architecture.
\\
\rowsep

\grouprow{Renewal and comprehensive reality}

\textbf{Non-repetition of tajalli}
&
Self-disclosure is not repeated in precisely
the same manner; creation is continuously renewed.
&
Thesis
&
Argued/Derived Claims
&
Non-repetition of manifestations does not
establish the loss or plurality of Identity.
\\
\rowsep

\textbf{Perfect Human}
&
A comprehensive human reality related to
the divine and cosmic dimensions.
&
Thesis
&
Conditionally Supported Claims
&
The concept has no direct equivalent in the
foundational ontology of Book I.
\\
\rowsep

\textbf{Five Divine Presences}
&
A systematic scheme developed in the later
Akbarian tradition, especially through Qunawi
and subsequent interpreters.
&
Historical interpretation
&
Conditionally Supported Claims
&
The later scheme must not be attributed
unqualifiedly to Ibn Arabi.
\\
\rowsep

\grouprow{Comparative status}

\textbf{Being and Identity}
&
Book I maintains Analytical Distinction,
Ontological Inseparability, and Absence
of Ontological Priority.
&
Foundational Commitment
&
Foundationally Adopted Commitments
&
These commitments are not derived from
Ibn Arabi's account.
\\
\rowsep

\textbf{Identity as fixed entity}
&
The identification of Book I Identity with
the Akbarian fixed entity is rejected.
&
Rejected Thesis
&
Rejected/Superseded Claims
&
The two concepts perform different
ontological and explanatory functions.
\\
\rowsep

\textbf{Comparative synthesis}
&
The two systems address overlapping questions
through distinct metaphysical architectures.
&
Thesis
&
Provisional Constructions
&
Conceptual resemblance does not establish
equivalence, derivation, or historical influence.
\\
\rowsep

\textbf{Plurality of identities}
&
The comparison does not settle whether
multiple identities obtain and under
which conditions.
&
Open Problem
&
Unresolved/Open Problems
&
No conclusion follows merely from the
plurality of names or manifestations.
\\

\end{longtable}

\endgroup

% =============================================================================
% END OF TABLE TL01
% =============================================================================



% =============================================================================
% SS19 — TRANSITION TO MULLA SADRA
% =============================================================================

\subsection{Transition to Mulla Sadra}
\label{subsec:ch08_s16_ss19}

% Structural ID: VS.B01.P01.C08.S16.SS19.P001
\phantomsection
\label{para:VS.B01.P01.C08.S16.SS19.P001}

The examination of Ibn Arabi has centered
on the unity of the Real, the determination
of possible things, the divine names,
self-disclosure, and the ontological
status of intermediate realities. His
account provides a theological metaphysics
in which the multiplicity of the cosmos
is related to one unrestricted source
of existence. The next comparison moves
toward a philosophical system that
develops the primacy and gradation of
existence and gives substantial motion
a central role in explaining the
dynamism of the existent.

% Structural ID: VS.B01.P01.C08.S16.SS19.P002
\phantomsection
\label{para:VS.B01.P01.C08.SS19.P002}

The transition to Mulla Sadra does not
establish that his philosophy is simply
a continuation of Ibn Arabi's system.
It introduces a distinct set of questions
concerning the primacy of existence,
the gradation of being, substantial
motion, and the relation between
existence and quiddity. The comparison
must therefore proceed by reconstructing
Sadra's own arguments and determining
their relation to the foundational
distinctions of Book I.

% =============================================================================
% END OF SECTION 16 OF CHAPTER 08
% =============================================================================
